
\chapter{The level-\(n\) protected scalar shadow and the \(K3\times E\) realization question}
\label{ch:zbps-realization}

The denominator algebra and the protected Pfaffian sit at the
primitive and chart levels of the Beilinson chain
\(\mathsf{P}\to\mathsf{C}\to\mathsf{S}\to\mathsf{Z}\to\mathsf{A}\).
After taking the protected trace, the same \(\Delta_5\) survives,
squared, as the partition function on the closed
\(K3\times E\to\mathrm{pt}\) cobordism, equivalently the level-\(n\)
modular-functor scalar trace of the realization formalism of Vol IV
\cite{ChiralBarCobarVolIV}.  The scalar shadow is exactly the object
allowed by the Beilinson cut: an automorphic identity,
not an algebra; a trace, not an operator; an orientation-forgetful
invariant, not an orientation lift.  The orientation-level object is
\(\mathfrak D_X^{\mathrm{DI}}\) of Chapter~5; the scalar identity is
the level-\(n\) shadow it casts.

Three distinct meanings of \(\Delta_5\) coexist and are tagged on
first occurrence: the automorphic section
\(\mathcal D_X=\Delta_5\) (View I --- Borcherds--Gritsenko--Nikulin); the
Borcherds--Kac--Moody denominator
\(\operatorname{den}(\mathfrak g_{\Delta_5})=64^{-1}\Delta_5(2Z)\)
(View II --- Gritsenko--Nikulin); the protected Pfaffian
\(\mathrm{Pf}_{\mathrm{prot}}(\mathfrak D_X^{\mathrm{DI}})\) (View
III --- Chapter~5).  The handle ``\(\Delta_5\)'' alone never
substitutes for any of the three.

\section{The level-\(n\) protected scalar shadow on \(K3\times E\to\mathrm{pt}\)}
\label{sec:zbps-scalar-shadow}

\subsection{The Oberdieck--Pixton scalar identity}

Let \(X=S\times E\), with \(S\) a smooth complex K3 surface and \(E\) a
smooth elliptic curve.  The reduced Donaldson--Thomas / stable-pair /
Gromov--Witten correspondence on \(X\), in the primitive K3-class
chamber, gives the normalized scalar representative
\[
Z^X_{\mathrm{OP/DT/PT}}=-4096\,\Delta_5^{-2}
\]
on the genus-\(2\) Siegel upper half-space \(\mathfrak H_2\).  The
unscaled level-\(n\) trace retained below is
\(Z_{\mathrm{BPS}}^{K3\times E}=\Delta_5^{-2}\).

\begin{convention}[OP chamber, sign branch, primitive class]
\label{conv:op-chamber}
The variables
\((T_{\mathrm{OP}},Q_{\mathrm{OP}},P_{\mathrm{OP}})\) are the
Oberdieck--Pixton variables of \cite[\S0]{OB}:
\(T_{\mathrm{OP}}\) is the K3-surface degree,
\(Q_{\mathrm{OP}}\) the K3-curve class, and
\(P_{\mathrm{OP}}\) the elliptic-genus variable
(equivalently \(P_{\mathrm{OP}}=-p_{\mathrm{DT}}\), giving the
\((-p_{\mathrm{DT}})^n\) DT-chamber).  The K3 curve class
\(\beta_h\) is taken \emph{primitive} of self-intersection
\(\beta_h\cdot\beta_h=2h-2\) (\cite{MaulikPandharipande}); the
identity \eqref{eq:zbps-op-dt} is the primitive-class branch.
The leading \(-1\) of \(Z^X_{\mathrm{OP}}=-(\chi_{10}^{\mathrm{OP}})^{-1}\)
is the OP scalar convention; the squared theta-leading
\(4096=64^2\) is the normalization conversion
\(\chi_{10}^{\mathrm{OP}}=4096^{-1}\Delta_5^2\).  None of these is the
orientation character.
\end{convention}

The scalar identity \eqref{eq:zbps-op-dt} is an identity in the OP chamber.  If the
retained compact Hall source lies in a different stability chamber,
comparison with \eqref{eq:zbps-op-dt} requires the chamber-compatible
Hall descent datum of Definition~\ref{def:chamber-compatible-hall-descent};
the normal-ordered lattice map alone does not transport Hall brackets.

Under Convention~\ref{conv:op-chamber} and \cite[Theorem~1]{OB},
\cite[Proposition~5]{OPand}, \cite{MaulikPandharipande},
\begin{equation}
\label{eq:zbps-op-dt}
Z^X_{\mathrm{DT}}
\;=\;
Z^X_{\mathrm{OP}}(P_{\mathrm{OP}},Q_{\mathrm{OP}},T_{\mathrm{OP}})
\;=\;
-\,4096\,\Delta_5(Z)^{-2},
\qquad
Z=\Omega\in\mathfrak H_2.
\end{equation}
The constant \(4096=64^2\) is the squared theta-leading coefficient
\([q^{1/2}r^{1/2}s^{1/2}]\Delta_5\) of the automorphic section, and
the leading sign is the Oberdieck--Pixton scalar convention
\(Z^X_{\mathrm{OP}}=-(\chi_{10}^{\mathrm{OP}})^{-1}\).  Both are
normalization data of the imported reduced theory.

\begin{theorem}[Oberdieck--Pixton scalar branch and unscaled Igusa shadow]
\label{thm:zbps-level-n-shadow}
\ClaimStatusProvedElsewhere
With variables
\((T_{\mathrm{OP}},Q_{\mathrm{OP}},P_{\mathrm{OP}})\) normalized in the
Oberdieck--Pixton chamber of the Oberdieck--Pixton normalization
\cite[Theorem~1]{OB}, the proved reduced DT/PT scalar branch is
\[
Z^X_{\mathrm{DT}}=Z^X_{\mathrm{OP}}=Z^X_{\mathrm{PT}}
=-4096\,\Delta_5(Z)^{-2}.
\]
\ClaimStatusDefinition\ In the unscaled trace convention of this
manuscript, the level-\(n\) protected scalar shadow on the closed
\(K3\times E\to\mathrm{pt}\) cobordism is defined by
\begin{equation}
\label{eq:zbps-equality}
Z_{\mathrm{BPS}}^{K3\times E}(T_{\mathrm{OP}},Q_{\mathrm{OP}},P_{\mathrm{OP}})
\;=\;
\bigl(\Phi_{10}^{\mathrm{un}}(Z)\bigr)^{-1}
\;=\;
\Delta_5(Z)^{-2}.
\end{equation}
Equivalently, in Oberdieck--Pixton normalization
\(\chi_{10}^{\mathrm{OP}}=2^{-12}\Delta_5^2\),
\(\Phi_{10}^{\mathrm{un}}=\chi_{10}=\Delta_5^2\), and the unscaled
form \((\Phi_{10}^{\mathrm{un}})^{-1}\) is the retained level-\(n\)
trace convention on each primitive closed K3-class fibre in the OP
chamber.
The normalization packet
\(\texttt{certificates/normalizations/delta5\_theta\_leading}\)
checks the named convention rows
\(\Delta_{10}\), \(\chi_{10}\), \(\Phi_{10}^{\mathrm{un}}\),
\(\chi_{10}^{\mathrm{OP}}\), \((\Phi_{10}^{\mathrm{un}})^{-1}\),
\((\chi_{10}^{\mathrm{OP}})^{-1}\), and \(Z^X_{\mathrm{OP}}\)
separately; the scalar equalities below are relations among those rows,
not a collapse of the conventions.
The divisor packet
\(\texttt{certificates/automorphic/delta5\_humbert\_divisor}\)
checks that the polar divisor of the inverse square
\(\Delta_5^{-2}\) has order \(2\) on the same Humbert support on which
\(\Delta_5\) has order \(1\); this is automorphic divisor arithmetic,
not the mirror-discriminant assertion of Chapter~\ref{ch:view4-mirror}.
The Behrend-weighted reduced quotient
integration of Oberdieck's reduced PT/DT scalar-quotient integration
\cite{OberdieckReducedPT} gives the OP chamber branch
\[
Z^X_{\mathrm{DT}}=Z^X_{\mathrm{OP}}=Z^X_{\mathrm{PT}}
=-4096\,Z_{\mathrm{BPS}}^{K3\times E}.
\]
Thus the OP branch is a normalized scalar representative of the same
unscaled trace, not the definition of \(Z_{\mathrm{BPS}}^{K3\times E}\).
\end{theorem}

\begin{proof}
Equation~\eqref{eq:zbps-equality} is the dictionary between the
Oberdieck--Pixton primitive series and the unscaled Igusa form
\(\Phi_{10}^{\mathrm{un}}\): the Oberdieck--Pixton normalization
\(\chi_{10}^{\mathrm{OP}}=D_5^2=64^{-2}\Delta_5^2=2^{-12}\Delta_5^2\)
of \cite[\S0]{OB} together with the Igusa identification
\(\Phi_{10}^{\mathrm{un}}=\chi_{10}=\Delta_5^2\) of
\cite{Igusa1962,Igusa1964II} gives
\((\Phi_{10}^{\mathrm{un}})^{-1}=\Delta_5^{-2}\).  Multiplying by
\(2^{12}=4096\) and the OP scalar sign restores
\eqref{eq:zbps-op-dt}.  The reduced DT/PT comparison is
\cite[Theorem~3(i)]{OberdieckReducedPT}, the reduced PT/GW comparison
on primitive K3 classes is
\cite[Proposition~5]{OPand}, and the OP product identity is
\cite[Theorem~1]{OB}; their composition gives the displayed equality
in the OP chamber.
\end{proof}

\subsection{Three readings of the same exponent}

The integers \(f(n,l)\) of the weight-zero index-one weak Jacobi form
\(\phi_{0,1}\) drive both sides of the OP/Borcherds dictionary.  At
the primitive datum they are virtual \(K_0\)-superdimensions of
\(\mathbb U_{n,l}\), and after taking scalar shadows they appear twice:
once as exponents of the Borcherds product on the automorphic side,
once as multiplicities in the Mukai dimension count of the
Hilbert-scheme stratification on the geometric side.  In
\(\mathfrak g_{\Delta_5}\), the same \(f(nm,l)\) is the signed
root supermultiplicity at \(\alpha(n,l,m)=2nf_2-lf_3+2mf_{-2}\).

The variable map of the Oberdieck--Pixton normalization \cite[Theorem~1]{OB} encodes
this triple reading.  A factor \((1-q^nr^ls^m)^{f(nm,l)}\) of the
Borcherds product on the automorphic side is, on the Oberdieck--Pixton
side, a factor
\((1-p_{\mathrm{OP}}^kq_{\mathrm{OP}}^h\widetilde q_{\mathrm{OP}}^d)^{c(4hd-k^2)}\)
of \(\chi_{10}^{\mathrm{OP}}\), the K3-elliptic-genus exponents
\(c(4hd-k^2)=2f(hd,k)\) recording the squared automorphic line.  The
multiplicities \(2f\) are exactly the K3 elliptic-genus coefficients
of \(Z_{\mathrm{K3}}=2\phi_{0,1}\); the squaring is part of the
identification \(\chi_{10}^{\mathrm{OP}}=D_5^2\) of
the Oberdieck--Pixton/Pandharipande scalar identity \cite[Theorem~1]{OB}.

\begin{remark}[Constants are normalization, not orientation]
\label{rem:zbps-constants-not-orientation}
Three constants intervene in \eqref{eq:zbps-op-dt}.  The leading
\(64=[q^{1/2}r^{1/2}s^{1/2}]\Delta_5\) is the theta-constant
normalization of the Igusa section; the monic Borcherds product is
\(D_5=64^{-1}\Delta_5\).  The factor \(4096=64^2\) is the
squared theta-leading coefficient produced by the OP convention
\(\chi_{10}^{\mathrm{OP}}=D_5^2\).  The leading minus sign is the OP
scalar branch \(Z^X_{\mathrm{OP}}=-(\chi_{10}^{\mathrm{OP}})^{-1}\).
None of the three is the orientation character.  The protected
orientation character \(\epsilon_o\) and the automorphic reflection
character \(\nu_{\Delta_5}\) are determined elsewhere: \(\nu_{\Delta_5}\)
by the Maass transformation law of \(\Delta_5\), \(\epsilon_o\) by the
rank-one type-II Pfaffian signs together with the retained
half-Hilbert wall atlas.  The number \(4096=2^{12}\) occurs in
two different structures, not as one orientation datum: the
Pfaffian-side \(2^{12}\) is the cardinality of the retained sign
orbit, while the OP-side \(2^{12}\) is
the squared theta-leading coefficient in the imported scalar identity.
\end{remark}

\subsection{Worked Fourier-coefficient check}

The leading Fourier coefficients of \(\Delta_5\) and \(\Delta_5^2\)
provide the elementary verification that the constants in
\eqref{eq:zbps-op-dt} are normalization, not enumerative content.
Recall \(\phi_{0,1}\) has leading expansion
\begin{equation}
\label{eq:zbps-phi01-leading}
\phi_{0,1}(\tau,z)
\;=\;
(r^{-1}+10+r)+q\bigl(10r^{-2}-64r^{-1}+108-64r+10r^2\bigr)+O(q^2),
\end{equation}
with coefficients
\(f(0,-1)=f(0,1)=1\), \(f(0,0)=10\), \(f(1,1)=f(1,-1)=-64\),
\(f(1,0)=108\), \(f(1,2)=f(1,-2)=10\), \(f(1,3)=f(1,-3)=0\).
The Borcherds--Gritsenko--Nikulin product
\begin{equation}
\label{eq:zbps-bgn-product}
\Delta_5(Z)
\;=\;
64\,q^{1/2}r^{1/2}s^{1/2}
\prod_{(n,l,m)\in\Gamma_{\mathrm{eff}}}
(1-q^nr^ls^m)^{f(nm,l)}
\end{equation}
opens with the theta-leading term
\([q^{1/2}r^{1/2}s^{1/2}]\Delta_5=64\); the next non-trivial
timelike coefficient belongs to the monic product after the leading
monomial is removed:
\[
[qrs]\,D_5/(qrs)^{1/2}=93,\qquad D_5=64^{-1}\Delta_5 .
\]
The isolated factor \((1-qrs)^{f(1,1)}=(1-qrs)^{-64}\) contributes
\(64qrs\), and the remaining root factors contribute the other
\(29qrs\), as in Appendix~\ref{app:borcherds-product-expansion}.
Squaring \eqref{eq:zbps-bgn-product},
\begin{equation}
\label{eq:zbps-delta5-squared-leading}
\Delta_5^2(Z)
\;=\;
4096\,q\,r\,s
\prod(1-q^nr^ls^m)^{2f(nm,l)},
\end{equation}
with theta-leading \(64^2=4096\).  The OP product
\(\chi_{10}^{\mathrm{OP}}=p\,q\,\widetilde q\prod(1-p^kq^h\widetilde
q^d)^{c(4hd-k^2)}\), with K3-elliptic exponents \(c(4hd-k^2)=2f(hd,k)\),
is exactly the right-hand side of \eqref{eq:zbps-delta5-squared-leading}
divided by \(4096\).  Thus
\begin{equation}
\label{eq:zbps-leading-coeff-check}
\chi_{10}^{\mathrm{OP}}=4096^{-1}\,\Delta_5^2,
\qquad
[pq\widetilde q]\,\chi_{10}^{\mathrm{OP}}=1,
\qquad
[q^{1/2}r^{1/2}s^{1/2}]\,(64^{-1}\Delta_5)=1.
\end{equation}
The squared theta-leading coefficient \(4096\) is therefore the
single explicit normalization constant in the OP/DT scalar
identity \(Z^X_{\mathrm{OP}}=-(\chi_{10}^{\mathrm{OP}})^{-1}\); the
sign \(-1\) is OP convention; together they produce the constant
\(-4096\) of \eqref{eq:zbps-op-dt}.  The Fourier coefficient at
\(q^{1/2}r^{1/2}s^{1/2}\) identifies the scalar \(4096\) as
\(\bigl([q^{1/2}r^{1/2}s^{1/2}]\Delta_5\bigr)^2\), a leading-coefficient
square, not an enumerative invariant.

\subsection{Cobordism positioning and the closed \(K3\times E\to\mathrm{pt}\)
cobordism}

The cobordism \(K3\times E\to\mathrm{pt}\) is the closed
Calabi--Yau-threefold target viewed in the retained trace convention.
The trace is the Laurent expansion of the
meromorphic automorphic section \(\Delta_5^{-2}\); evaluation at a
point of \(\mathfrak H_2\) outside the polar divisor gives a complex
number.  The ``level-\(n\)'' label is the curve-counting integer
\[
n=\chi(\mathcal O_Y)
\]
for a one-dimensional subscheme \(Y\subset K3\times E\) of class
\((\beta_h,d)\).  The target has
\[
\chi(\mathcal O_{K3\times E})
=\chi(\mathcal O_{K3})\chi(\mathcal O_E)=2\cdot 0=0,
\]
and this target Euler characteristic is not the exponent \(n\).  The
Laurent variables record \(n\), the elliptic degree \(d\), and the K3
norm \(h-1\).  The dual notation
\(Z_{\mathrm{BPS}}^{K3\times E}(T_{\mathrm{OP}},Q_{\mathrm{OP}},P_{\mathrm{OP}})\)
stores the elliptic and K3
gradings as a Laurent expansion of the inverse Igusa cusp form on
\(\mathfrak H_2\).

The \(K3\times E\) target is a threefold; the perfect obstruction
theories and reduced virtual fundamental classes live on the
stable-pair and DT moduli spaces
\(P_n(X,(\beta_h,d))\) and \(I_n(X,(\beta_h,d))\), not on the target
itself and not on any auxiliary fourfold.  The
closed K3-class fibre at primitive curve class \(\beta_h\) of norm
\(2h-2\) supplies the K3 component of the coefficient indexed by
\((h,n,d)\), not a new target dimension.

\section{Spin \(L\)-factor normalization for \(\Delta_5^2\)}
\label{sec:spin-l-factor}

\subsection{The Saito--Kurokawa eigenline}

The unscaled Igusa form \(\Phi_{10}^{\mathrm{un}}=\Delta_5^2\) is the
weight-\(10\) Saito--Kurokawa lift of the unique normalized cusp form
\(g=\Delta E_6\in S_{18}(\SL_2(\Z))\); the Hecke eigenline is the
one-dimensional \(\C\Delta_{10}^{\theta}=\C\chi_{10}^{\mathrm{OP}}\),
independent of the scalar normalization of its generator.

\begin{proposition}[Hecke eigenvalues on the Saito--Kurokawa
eigenline]
\label{prop:zbps-sk-hecke}
For each prime \(p\), let \(T_1(p)\) and \(T_2(p^2)\) denote the
degree-one and degree-two genus-two Hecke operators on
\(S_{10}(\Sp_4(\Z))\) of \cite[\S 4]{Andr74}.  The Saito--Kurokawa
form \(\Delta_{10}^{\theta}=\Delta_5^2\) is a simultaneous
eigenform with eigenvalues
\[
 T_1(p)\,\Delta_{10}^{\theta}
 \;=\;
 \lambda_1^{\mathrm{SK}}(p)\,\Delta_{10}^{\theta},
 \qquad
 T_2(p^2)\,\Delta_{10}^{\theta}
 \;=\;
 \lambda_2^{\mathrm{SK}}(p^2)\,\Delta_{10}^{\theta},
\]
where the eigenvalues are determined by the Saito--Kurokawa
correspondence from the elliptic Hecke eigenvalues of
\(g=\Delta E_6\):
\[
 \lambda_1^{\mathrm{SK}}(p)\;=\;\lambda_g(p)+p^8+p^9,
 \qquad
 \lambda_2^{\mathrm{SK}}(p^2)\;=\;\lambda_g(p)^2+\lambda_g(p)\bigl(p^8+p^9\bigr)+p^{18},
\]
read off the \(X\)- and \(X^2\)-coefficients of the local Euler factor
associated with \eqref{eq:zbps-spinor-factor}, with \(\lambda_g(p)\) the normalized \(p\)-th Hecke eigenvalue of
\(g=\Delta E_6\), the unique normalized cusp form in
\(S_{18}(\SL_2(\Z))\): from \(\Delta E_6 = q-528q^2-4284q^3+\cdots\),
\(\lambda_g(2)=-528\), \(\lambda_g(3)=-4284\) \cite[Table~1]{ZagierSK}.
Concretely \(\lambda_1^{\mathrm{SK}}(2)=-528+2^8+2^9=240\) and
\(\lambda_1^{\mathrm{SK}}(3)=-4284+3^8+3^9=21960\), while
\(\lambda_2^{\mathrm{SK}}(4)=135424\) and
\(\lambda_2^{\mathrm{SK}}(9)=293343849\); the
Ramanujan bound is violated (\(\lvert\lambda_1^{\mathrm{SK}}(p)\rvert
\sim p^9\), whereas tempered classical eigenvalues at weight \(10\)
have size \(p^{17/2}\)), the signature of the
Saito--Kurokawa CAP form.
\end{proposition}

\begin{proof}
The Saito--Kurokawa correspondence \cite[\S 1]{ZagierSK} produces, for
each weight \(2k\) elliptic newform \(g\in S_{2k}(\SL_2(\Z))\), a
weight-\(k+1\) Siegel cusp form \(F_g\in S_{k+1}(\Sp_4(\Z))\) with
spinor \(L\)-function
\(L_{\mathrm{spin}}(s,F_g)=\zeta(s-k+1)\zeta(s-k)L(s,g)\) (which is
exactly \eqref{eq:zbps-spinor-factor} at \(k=9\)).  The Euler product
factorisation determines the Hecke eigenvalues at each prime.  In
Andrianov's polynomial the Siegel weight is \(K=k+1\), so
\(p^{2K-4}=p^{2k-2}\), \(p^{2K-3}=p^{2k-1}\), and
\(p^{4K-6}=p^{4k-2}\).  Writing
\(L_{\mathrm{spin}}(s,F_g)=\prod_p L_p(p^{-s},F_g)^{-1}\) with
\[
 L_p(X,F_g)
 \;=\;(1-p^{k-1}X)(1-p^kX)(1-\lambda_g(p)X+p^{2k-1}X^2)
\]
and comparing with the Andrianov factorisation
\cite[\S 4.5]{Andr74}
\[
L_p(X,F_g)=
1-\lambda_1^{\mathrm{SK}}(p)X+
(\lambda_1^{\mathrm{SK}}(p)^2-\lambda_2^{\mathrm{SK}}(p^2)-p^{2k-2})X^2-
\lambda_1^{\mathrm{SK}}(p)p^{2k-1}X^3+p^{4k-2}X^4,
\]
the \(X\)-coefficient gives
\(\lambda_1^{\mathrm{SK}}(p)=\lambda_g(p)+p^{k-1}+p^k\).  The
\(X^2\)-coefficient of the Euler product is
\[
2p^{2k-1}+(p^{k-1}+p^k)\lambda_g(p).
\]
Since the Andrianov coefficient of \(X^2\) is
\(\lambda_1^{\mathrm{SK}}(p)^2-\lambda_2^{\mathrm{SK}}(p^2)-p^{2k-2}\),
one obtains
\[
\begin{aligned}
\lambda_2^{\mathrm{SK}}(p^2)
&=(\lambda_g(p)+p^{k-1}+p^k)^2-p^{2k-2}
  -2p^{2k-1}-(p^{k-1}+p^k)\lambda_g(p) \\
&=\lambda_g(p)^2+\lambda_g(p)(p^{k-1}+p^k)+p^{2k}.
\end{aligned}
\]
At \(k=9\) this is the displayed formula.  The dimension formula
\(\dim S_{18}(\SL_2(\Z))=1\) forces \(g=\Delta E_6\) at \(k=9\), and
the eigenvalues at \(p=2,3\) are tabulated in
\cite[Table~1]{ZagierSK}.
\end{proof}

\begin{remark}[The SK-correspondence onto the full-level Maa{\ss} line]
\label{rem:zbps-sk-hecke-level0}
The SK-correspondence \(g\mapsto F_g\) is a Hecke-equivariant
injection \(S_{2k}(\SL_2(\Z))\hookrightarrow S_{k+1}(\Sp_4(\Z))\)
whose image is the Maa{\ss} subspace
\(\mathrm{Maass}_{k+1}\subset S_{k+1}(\Sp_4(\Z))\).  At \(k=9\),
\(\dim\mathrm{Maass}_{10}=\dim S_{18}(\SL_2(\Z))=1\) and
\(\mathrm{Maass}_{10}=\C\Delta_5^2\); the Hecke eigenvalues
\(\lambda_i^{\mathrm{SK}}\) of
Proposition~\ref{prop:zbps-sk-hecke} record the explicit
arithmetic translation of the full-level elliptic Hecke action on
\(g=\Delta E_6\) under the SK-correspondence.  No new modular
forms or arithmetic data enter; the Saito--Kurokawa eigenline of
\(\Delta_5^2\) is fully determined by the full-level arithmetic of
\(\Delta E_6\) and the SK-correspondence.
\end{remark}
The OP
monic normalization is
\begin{equation}
\label{eq:zbps-op-monic}
\chi_{10}^{\mathrm{OP}}
\;=\;
2^{-12}\,\Delta_{10}^{\theta}
\;=\;
4096^{-1}\,\Delta_5^2,
\end{equation}
and the level-\(n\) shadow~\eqref{eq:zbps-equality} is the inverse of
the \emph{theta-normalized} Hecke eigenline element
\(\Delta_{10}^{\theta}=\Delta_5^2\).

\begin{theorem}[Andrianov factorization of the spinor \(L\)-function]
\label{thm:zbps-andrianov-factor}
\ClaimStatusProvedElsewhere
For the Saito--Kurokawa representation generated by
\(\Delta_{10}^{\theta}=\Delta_5^2\),
\begin{equation}
\label{eq:zbps-spinor-factor}
L_{\mathrm{spin}}(s,\Delta_{10}^{\theta})
\;=\;
\zeta(s-8)\,\zeta(s-9)\,L(s,g),
\qquad g=\Delta E_6.
\end{equation}
Equivalently, for Andrianov's local polynomial
\[
Q_{p,F}(t)=
1-\lambda_F(p)t+
\bigl(\lambda_F(p)^2-\lambda_F(p^2)-p^{16}\bigr)t^2
-\lambda_F(p)p^{17}t^3+p^{34}t^4
\]
at weight \(10\), the Euler product
\(\prod_p Q_{p,F}(p^{-s})^{-1}\) is the left-hand side of
\eqref{eq:zbps-spinor-factor}.  Andrianov's completed function is
\[
\Psi_F(s)
=(2\pi)^{-s}\Gamma(s)\Gamma(s-8)
L_{\mathrm{spin}}(s,F).
\]
\end{theorem}

This is Andrianov's spin \(L\)-factor factorization
\cite[Theorem~3.1.1]{Andr74}; the polynomial \(Q_{p,F}\) is
Andrianov's genus-two Hecke polynomial in weight \(10\), with
coefficient
\(\lambda_F(p)^2-\lambda_F(p^2)-p^{2\cdot 10-4}\) at \(t^2\).
The proof uses Andrianov's spinor \(L\)-factor formula for the
degree-\(2\) Saito--Kurokawa lift attached to the weight-\(18\)
elliptic cusp form \(g\)
\cite{Andr74}, together with the dimension count
\(\dim S_{18}(\SL_2(\Z))=1\) which forces the source form to be
\(g=\Delta E_6\) \cite{ZagierSK}.

\begin{proposition}[Polar coefficient at \(s=10\)]
\label{prop:zbps-spin-pole-residue}
\ClaimStatusProvedElsewhere
With \(\Lambda(g,s)=(2\pi)^{-s}\Gamma(s)L(s,g)\), the functional
equation \(\Lambda(g,s)=-\Lambda(g,18-s)\) gives \(L(9,g)=0\), so the
elliptic central zero cancels the possible \(\zeta(s-8)\)-pole in the
spin product at \(s=9\).  At \(s=10\) the polar coefficient is
\begin{equation}
\label{eq:zbps-residue}
\operatorname*{Res}_{s=10}L_{\mathrm{spin}}(s,\Delta_{10}^{\theta})
\;=\;
\zeta(2)\,L(10,g)
\;=\;
\frac{\pi^2}{6}\,L(10,g),
\end{equation}
with
\[
L(10,g)/(2\pi i)^{10}\Omega^+(g)\in\Q
\]
by Manin--Deligne periods.  The period theorem fixes the
normalisation of the critical value; it is not a non-vanishing theorem.
Thus \(L_{\mathrm{spin}}(s,\Delta_{10}^{\theta})\) has a simple pole at
\(s=10\) exactly when \(L(10,g)\ne0\).
\end{proposition}

This is the Saito--Kurokawa critical-value identity \cite{ZagierSK} together with
the Manin--Deligne period decomposition \cite{Manin72,Del79}; the proof uses the
Eichler--Shimura--Manin period theorem for the rational newform \(g\)
\cite{Manin72,KZ}, in Deligne's period normalization \cite{Del79}.

\begin{remark}[Hecke eigenline is normalization-independent]
The spinor \(L\)-function depends on the Hecke eigenline
\(\C\Delta_{10}^{\theta}=\C\chi_{10}^{\mathrm{OP}}\), not on the scalar
used to label its generator.  Theorem~\ref{thm:zbps-andrianov-factor}
and Proposition~\ref{prop:zbps-spin-pole-residue} are statements about
the eigenline.  The OP scalar constant \(2^{-12}\) of
\eqref{eq:zbps-op-monic} relabels the generator from
\(\Delta_5^2\) to \(\chi_{10}^{\mathrm{OP}}\); it does not change
\(L_{\mathrm{spin}}\), the residue, or the cancellation of the
elliptic central zero against the zeta pole at \(s=9\).
\end{remark}

\subsection{Independence from the Dirac--Igusa construction}

The Saito--Kurokawa factorization of \(L_{\mathrm{spin}}(s,\Delta_5^2)\)
and the OP/DT scalar identity~\eqref{eq:zbps-op-dt} are imported
results: the first from Andrianov \cite{Andr74}, the second from
Oberdieck--Pixton \cite{OB} and Oberdieck \cite{OberdieckReducedPT}.
Neither uses the protected Pfaffian \(\mathrm{Pf}_{\mathrm{prot}}\) of
Chapter~5, the orientation character \(\epsilon_o\) of (O1)/(O1\(^+\)),
the type-II wall normal form of (O2), or the primitive recognition
\(P_X\cong\mathfrak g_{\Delta_5}\) problem of Chapter~6.  Conversely, the
construction of \(\mathfrak D_X^{\mathrm{DI}}\) and the
Pfaffian--Dirac theorem \(\mathrm{Pf}_{\mathrm{prot}}(\mathfrak
D_X^{\mathrm{DI}})=\Delta_5\) do not use the spinor \(L\)-factor or
the scalar inverse Igusa form.  The two routes meet at the squared
section \(\Delta_5^2=\Phi_{10}^{\mathrm{un}}\) and at no earlier
point.

\section{What the scalar shadow does not supply}
\label{sec:scalar-trace-not-gravity}

\subsection{\texorpdfstring{$Z_{\mathrm{BPS}}$ is the level-$n$ protected trace, not a path integral}{Z\_BPS is the level-n protected trace, not a path integral}}

The squared section \(\Delta_5^2=\Phi_{10}^{\mathrm{un}}\) is the
determinant section, and its inverse
\(Z_{\mathrm{BPS}}^{K3\times E}=\Delta_5^{-2}\) is the protected
scalar trace in the Beilinson chain, after the
primitive object at level \(\mathsf{P}\) (the open factorization
dg-category, the boundary vacuum \(b\), and \(A_b=\mathrm{End}_{\mathcal
C}(b)\)) and after the chart object at level \(\mathsf{C}\)
(the disk algebra and its bar / cobar / chiral Hochschild data).  The
scalar partition function is the \emph{trace} of the protected
data; it does not by itself construct the level-\(\mathsf A\) acting
algebra.

At finite height the protected trace is a separate row-level datum:
\[
\mathrm{trace\_categories},\quad
\mathrm{protected\_trace\_functors},\quad
\mathrm{trace\_operators},\quad
\mathrm{scalar\_normalizations},\quad
\mathrm{trace\_identities},\quad
\mathrm{forgetful\_maps},\quad
\mathrm{gravity\_residuals},\quad
\mathrm{transitions},\quad
\mathrm{scalar\_firewall}.
\]
These rows record the level-\(\mathsf Z\) trace category, the closed
\(K3\times E\to\mathrm{pt}\) trace functor, the operator whose
determinant square is \(\Phi_{10}^{\mathrm{un}}\), the constants
\(64\), \(4096\), and the OP sign \(-1\), the identity
\(\Delta_5^{-2}=(\Phi_{10}^{\mathrm{un}})^{-1}\), and the forgetful
maps killing orientation, primitive bracket, Hopf pairing, PBW, and
parity data.  The gravity-residual rows must mark the Hall--Borcherds
level-\(\mathsf A\) operator algebra as conjectural, not constructed,
and unused in the trace proof.  The scalar firewall excludes
orientation characters, Pfaffian lines, primitive brackets, Hopf
pairings, PBW bases, parity splittings, source Koszul maps, and
gravitational path integrals as data recovered from the scalar trace.

\begin{scope}[No scalar trace is a 3D gravitational path integral]
\label{scope:zbps-no-3d-gravity}
\ClaimStatusOpen
The identity~\eqref{eq:zbps-op-dt}, proved through OP/DT/PT primary
literature, is a closed automorphic equation between scalar power
series.  Neither this scalar identity, nor the automorphic section of
Chapter~3, nor the protected Pfaffian of Chapter~5 identifies
\(Z_{\mathrm{BPS}}^{K3\times E}\) to a \(3\)-dimensional
gravitational path integral on \(\mathrm{AdS}_3\times S^3\times K3\),
to a metric Euclidean path-integral over the full conformal class of
the boundary, or to a holographic dual partition function in any sense
beyond the boundary-CFT reading recorded by Vol~II.  The equality
\(\Phi_{10}^{\mathrm{un}}=\Delta_5^2\)
gives the determinant section of a candidate gravity-line theory; the
partition character would be the reciprocal section
\(\Delta_5^{-2}\).  The construction of the gravity-line operator
algebra is open.
\end{scope}

\begin{remark}[Protected trace versus path integral]
The assertion that \(Z_{\mathrm{BPS}}\) is a \(3d\) gravitational
path integral identifies a protected scalar trace with a
level-\(\mathsf A\) acting theory.  The precise statement is that
\(Z_{\mathrm{BPS}}^{K3\times E}\) is the level-\(n\) protected trace of
the retained \(K3\times E\) data; the
gravitational path integral, if such a thing exists at this boundary,
requires the Hall--Borcherds residual of Vol~II
\textup{(\S\ref{sec:hall-borcherds-residual})}.  The residual is open;
the path-integral realization is not claimed.  In particular, no
microscopic state space \(\mathcal H_{K3\times E}\) is asserted to
support \(\dim\mathcal H_{K3\times E,\gamma}=f(\gamma)\) outside the
charge-graded virtual reading recorded in Chapter~2.
\end{remark}

\subsection{The Hall--Borcherds residual}
\label{sec:hall-borcherds-residual}

The passage from a scalar partition function to a chain-level
gravity-line operator algebra requires a comparison datum that lives
neither in the automorphic / OP reduced theory nor in the compact
Pfaffian datum.  Vol~II names this comparison.

\begin{conjecture}[Hall--Borcherds residual; gravity-line comparison]
\label{conj:zbps-hall-borcherds-residual}
\ClaimStatusConjectured
Let \(X=K3\times E\) and assume the oriented critical Hall datum used
for the Vol~III positive-half comparison
\begin{equation}
\label{eq:zbps-hall-to-borch}
\Theta_{\mathrm{Hall}\to\mathrm{Borch}}\colon
\widehat{\mathrm{CoHA}}^{\mathrm{red}}_{\Lambda^{2,1}_{II}}(X)
\longrightarrow
\widehat{U^{\mathrm{ch}}\bigl(\mathfrak n_+(\mathfrak g_{\Delta_5})\bigr)}.
\end{equation}
The Hall--Borcherds residual is the conjectural extension of
\eqref{eq:zbps-hall-to-borch}, after Drinfeld doubling, current
enveloping along \(E\), and weight-completion, to a filtered
\(\mathsf{SC}^{\mathrm{ch,top}}\) morphism from the \(K3\times E\) Hall
boundary object to the gravitational boundary line algebra, whose
derived-centre trace has scalar character
\((\Phi_{10}^{\mathrm{un}})^{-1}=\Delta_5^{-2}\).
\end{conjecture}

The residual is the Vol~II gravity-line Hall--Borcherds residual
\cite{ChiralBarCobarVolII}: it asks for the gravity-line operator
algebra acting on the boundary.  It is not claimed here; its existence
and the construction of its witnessing morphism are open.

\begin{remark}[Compatibility with the \(\kappa\)-stratum]
\label{rem:zbps-kappa-stratum}
The \(\kappa\)-tuple of the K3\(\times\)E witness consists of the structural data
behind the scalar shadow:
\begin{equation}
\label{eq:zbps-kappa-tuple}
\bigl(\kappa_{\mathrm{cat}},\,\kappa_{\mathrm{ch}}^{\mathrm{Hodge}},\,
\kappa_{\mathrm{ch}}^{\mathrm{Heis}},\,
\kappa_{\mathrm{BKM}},\,\kappa_{\mathrm{fiber}}\bigr)_{K3\times E}
\;=\;
(0,\,0,\,3,\,5,\,24).
\end{equation}
The fourth entry \(\kappa_{\mathrm{BKM}}=5\) is the cusp-form weight
\(\mathrm{wt}(\Delta_5)\), the \(\Delta_5\) row;
the fifth entry \(\kappa_{\mathrm{fiber}}=\chi_{\mathrm{top}}(K3)=24\)
is the topological Euler characteristic of the K3 fibre, distinct from
\(\chi(\mathcal O_{K3})=2\).  The additive formula
\(\kappa_{\mathrm{BKM}}=\kappa_{\mathrm{ch}}^{\mathrm{Hodge}}+\chi(\mathcal
O_{\mathrm{fiber}})\) confuses these two: with
\(\kappa_{\mathrm{ch}}^{\mathrm{Hodge}}=0\) and
\(\chi(\mathcal O_{K3})=2\), it predicts \(\kappa_{\mathrm{BKM}}=2\),
which fails because the universal Borcherds weight at \(N=1\) is
\(\kappa_{\mathrm{BKM}}(\Phi_1)=c_1(0)/2=10/2=5\) (Borcherds 1995,
Gritsenko 1999).  The five entries of the tuple are five \emph{distinct}
constructions; they cannot be combined additively.  The Vol~III
counterexample at \(N=1\) is the controlling instance.
\end{remark}

\subsection{Local-to-global in mixed holomorphic--topological language}

A second route from local protected data to a global compact
factorization algebra
runs through the mixed holomorphic--topological model of Costello and
the holomorphic de Rham obstruction class.  Formal local Hamiltonian BF
data imply such a global object only after the global holomorphic de
Rham obstruction class
\([\mathrm{obs}_{HT,\mathrm{glob}}]\in H^1(X,\mathcal O_X)\),
plus descent, quantum master equation, anomaly, and locality data,
with the Hamiltonian condition supplied by the local model.  For
\(X=K3\times E\) the global obstruction group is
\(H^1(K3\times E,\mathcal O_{K3\times E})\cong\C\) (non-zero), so the
formal-local-Hamiltonian-BF route does not by itself identify
\(Z_{\mathrm{BPS}}^{K3\times E}=\Delta_5^{-2}\) to a global compact
factorization algebra.  Both bridges --- Hall--Borcherds for the Vol~II gravity
line, mixed-HT for the Costello local-to-global ascent --- are open;
the only established object here is the scalar shadow.

\section{The \(K3\times E\) threefold realization question}
\label{sec:k3-e-threefold-realization}

The square root in the title is meant in Dirac's sense.  The OP/DT
scalar branch \(Z^X_{\mathrm{OP/DT}}=-4096\,\Delta_5^{-2}\) is
orientation-forgetful: the sign character is lost under squaring, and the
inverse-square scalar cannot recover the sign by which a Pfaffian-line
section distinguishes \(\Delta_5\) from \(-\Delta_5\).  The missing
first-order datum is the protected Pfaffian, whose square is the
inverse of the protected scalar shadow.

\subsection{From the scalar trace to the retained operator datum}

The Dirac square-root pattern compares the chiral-shadow Pfaffian with
the scalar trace.  Before taking the trace one has the protected
Pfaffian of View~\textcircled{3}; after the trace one has the inverse
square as scalar shadow.  The Pfaffian carries the orientation data;
inverting its square forgets it and gives the orientation-forgetful
scalar invariant.

\begin{theorem}[Orientation-forgetting scalar consequence of the Pfaffian]
\label{thm:zbps-square-root-consequence}
\ClaimStatusConditional
On \(K3\times E\), after the retained D0-HN datum, retained
quotient-orientation datum, chosen Weyl lifts, and retained
\((O2_{\mathrm{atlas}})\) wall datum, together with the retained finite
geometric Pfaffian datum \((P_{\mathrm{fin}})\), of
Appendix~\ref{app:obstruction-discharges}, the
protected Pfaffian \(\mathrm{Pf}_{\mathrm{prot}}(\mathfrak
D_X^{\mathrm{DI}})\) of Chapter~5 is the automorphic section
\(\mathcal D_X=\Delta_5\) of Chapter~3, and its inverse square is the
level-\(n\) protected scalar shadow:
\begin{equation}
\label{eq:zbps-square-root}
\bigl(\mathrm{Pf}_{\mathrm{prot}}(\mathfrak D_X^{\mathrm{DI}})\bigr)^{-2}
\;=\;
\Delta_5(Z)^{-2}
\;=\;
\bigl(\Phi_{10}^{\mathrm{un}}(Z)\bigr)^{-1}
\;=\;
Z_{\mathrm{BPS}}^{K3\times E}(T_{\mathrm{OP}},Q_{\mathrm{OP}},P_{\mathrm{OP}}),
\end{equation}
and the OP/DT scalar branch is
\(-4096\,Z_{\mathrm{BPS}}^{K3\times E}\).
\end{theorem}

\begin{proof}
The Pfaffian--Dirac theorem of Chapter~6, with
(D0), (O1), (O1\(^+\)), and \((O2_{\mathrm{atlas}})\) represented by
the retained data of Theorem~\ref{thm:G-four-obstruction-criterion}, and
with the retained finite geometric Pfaffian datum \((P_{\mathrm{fin}})\),
asserts
\(\mathrm{Pf}_{\mathrm{prot}}(\mathfrak D_X^{\mathrm{DI}})=\mathcal D_X=\Delta_5\).
Squaring,
\(\mathrm{Pf}_{\mathrm{prot}}(\mathfrak D_X^{\mathrm{DI}})^2
=\Delta_5^2=\Phi_{10}^{\mathrm{un}}\); inverting,
\(\mathrm{Pf}_{\mathrm{prot}}(\mathfrak D_X^{\mathrm{DI}})^{-2}
=(\Phi_{10}^{\mathrm{un}})^{-1}\).  By Theorem~\ref{thm:zbps-level-n-shadow}
the right-hand side equals the level-\(n\) protected scalar shadow.
The Oberdieck--Pixton/Pandharipande scalar identity
\cite[Theorem~1]{OB} says that the OP chamber representative is
\(-4096\) times this trace; the factor is the OP convention, the
squared theta-leading coefficient, and the sign branch of
\eqref{eq:zbps-op-dt}.
\end{proof}

\subsection{The orientation datum the inverse-square scalar cannot recover}

The square-root sign is invisible to the scalar.  An orientation-line
sign \(\varepsilon\in\{\pm 1\}\) acting by \(\Delta_5\mapsto\varepsilon\Delta_5\)
satisfies \(\varepsilon^2=1\); the scalar \(\Delta_5^{-2}\) is therefore
\(\varepsilon\)-invariant.  Selecting the sign requires a pre-trace
Pfaffian-orientation datum of View~\textcircled{3}: the protected orientation character
\(\epsilon_o\) of the type-II Pfaffian wall normal form (O2) and its
restriction to the type-II Weyl group \(W^{(2)}(\Lambda^{2,1}_{II})\),
matched against the automorphic reflection character \(\nu_{\Delta_5}\).
Equality
\begin{equation}
\label{eq:zbps-orientation-equality}
\epsilon_o\bigm|_{W^{(2)}(\Lambda^{2,1}_{II})}
\;=\;
\nu_{\Delta_5}\bigm|_{W^{(2)}(\Lambda^{2,1}_{II})}
\end{equation}
is the orientation-monodromy-character identity of
Chapter~\ref{ch:pfaffian-dirac}; Appendix~\ref{app:obstruction-discharges}
states the retained-data criteria for (D0), (O1), (O1\(^+\)), and (O2)
on \(K3\times E\).  No constant
appearing in \eqref{eq:zbps-op-dt} --- not the leading \(64\), not
the squared \(4096\), not the OP scalar minus sign --- supplies
\(\epsilon_o\).  The local sign is checked on a rank-one Pfaffian wall;
the global character requires the retained quotient orientation and
the chosen Weyl transport, not a normalization constant.

\subsection{The four obstructions and the realization question}

Let \(\mathfrak D_X^{\mathrm{DI}}\) denote the Dirac--Igusa pro-object
\begin{equation}
\label{eq:zbps-dirac-igusa-pro-object}
\mathfrak D_X^{\mathrm{DI}}
\;=\;
\lim_R\bigl(
\mathcal A^{E_3}_{X,R},\,F^{\mathrm{hyb}}_{X,R},\,\Gamma_{X,R},\,\Pi_X,\,
\Phi_R,\,o_R,\,H_R,\,P_R^\Pi,\,C_{X,R},\,\Theta_{\mathrm{Kos},R},\,
\mathcal L_{\mathrm{Pf},R},\,\mathrm{pf}_{X,R},\,\varepsilon_{o,R},\,
\mathrm{Rec}_R\bigr)
\end{equation}
of Chapter~5: the inverse limit, over finite root windows \(R\),
of fourteen pieces of data including the chosen finite
\(E_3\)-prefactorization algebra \(\mathcal A^{E_3}_{X,R}\), the hybrid
\(\operatorname{Ran}^{\mathrm{hyb}}(E)\) carrier
\(F^{\mathrm{hyb}}_{X,R}\), the cosection-twisted reduced d-critical
orientation \(o_R\), the Hall correspondence datum \(H_R\), the
Pfaffian line \(\mathcal L_{\mathrm{Pf},R}\), the Pfaffian section
\(\mathrm{pf}_{X,R}\), and the recognition map \(\mathrm{Rec}_R\).
The retained D0-HN datum, retained quotient-orientation datum, chosen
Weyl lifts, retained \((O2_{\mathrm{atlas}})\) datum, and retained finite
geometric Pfaffian datum \((P_{\mathrm{fin}})\) are the
Pfaffian and orientation data of Appendix~\ref{app:obstruction-discharges}
for
Theorem~\ref{thm:zbps-square-root-consequence}.

\begin{figure}[!ht]
\centering
\fbox{\parbox{0.94\textwidth}{
\textbf{Realization question (the K3\(\times\)E threefold) — retained scalar-trace level.}\;
The square-root identity~\eqref{eq:zbps-square-root} is a
theorem on \(K3\times E\) at the Pfaffian and scalar-trace level relative to the
retained D0-HN datum, retained quotient-orientation datum, chosen Weyl lifts,
retained \((O2_{\mathrm{atlas}})\) wall datum, and retained finite
geometric Pfaffian datum \((P_{\mathrm{fin}})\) in
Appendix~\ref{app:obstruction-discharges}:

\smallskip
\noindent
\((D0)\) \emph{D0-degeneration limit of the cosection-reduced
d-critical orientation theory.}
\hfill
\textbf{Criterion}: Theorem~\ref{thm:G-D0}, relative to the retained
D0-HN datum of Definition~\ref{def:G-D0-HN-system}: flat
D0-degeneration family, additive K3 semiregularity cosections,
vanishing-cycle and orientation specialization, Pfaffian-line
extension, protected-integration compatibility, and strict HN
transitions.

\smallskip
\noindent
\((O1)\) \emph{Strong reduced orientation on the
\(K3\times E\)-quotient self-Ext frame bundle.}
\hfill
\textbf{Relative to \((O1_{\mathrm{quot}})\)}:
Theorem~\ref{thm:G-O1}.

\smallskip
\noindent
\((O1^+)\) \emph{Weyl-equivariant transport of the orientation along
\(W^{(2)}(\Lambda^{2,1}_{II})\) reflections.}
\hfill
\textbf{Relative to the chosen Weyl lifts}:
Theorem~\ref{thm:G-O1-plus}.

\smallskip
\noindent
\((O2)\) \emph{Local Pfaffian wall normal form on type-II walls; the
\(2^{12}=4096\) sign sum on the half-Hilbert orbit, with
the \((R1)\)--\((R2)\) scalar separation.}
\hfill
\textbf{Relative to \((O2_{\mathrm{atlas}})\)}:
Theorems~\ref{thm:G-O2-orbit-transport},
\ref{thm:G-R1-R2}.

\smallskip
\noindent
\((P_{\mathrm{fin}})\) \emph{Finite geometric Pfaffian datum on a
cofinal HN system: skew perfect complexes, Pfaffian lines and sections,
type-II wall divisor orders, automorphic line comparisons, support and
parity rows, leading coefficient, and Pfaffian-line
Mittag--Leffler compatibility.}
\hfill
\textbf{Relative to \((P_{\mathrm{fin}})\)}:
Definition~\ref{def:G-finite-pfaffian-section}.

\smallskip
\noindent
The level-\(\mathsf Z\) Hall--Borcherds trace comparison is relative to
the retained \((Z_{\mathrm{HB}})\) datum and
Theorem~\ref{thm:G-vol2-discharge}: the retained trace datum places the
compact source in the level-\(\mathsf Z\) chiral-Hochschild domain, and
Vol~I Theorem~H \cite{ChiralBarCobarVolI} together with Vol~II
\cite[chiral-Hochschild Beilinson stratification]{ChiralBarCobarVolII}
clauses~(i) and~(ii) give the chain-level chiral-Hochschild
trace
\(\mathrm{Tr}^{\mathrm{bulk}}_n=\Delta_5^{-2}\) on
\(Z^{\mathrm{der}}_{\mathrm{ch}}(C_X)\).  The Z\(_{\mathrm{BPS}}\)
square-root identity is therefore a retained-datum theorem on
\(K3\times E\) at level \(\mathsf Z\), granted \((Z_{\mathrm{HB}})\)
and those two cited theorems; the
level-\(\mathsf A\) gravity-line operator algebra remains open in
Vol~II.  The finite packet
\(\texttt{certificates/trace/k3e\_protected\_trace}\) currently records
\((Z_{\mathrm{HB}})\) only by the obstruction ledger
\(\texttt{TRACE\_OBSTRUCTION\_LEDGER\_VERIFIED}\), with
\(\texttt{protected\_trace\_certification=false}\).}}
\end{figure}

\subsection{The threefold target, not a fourfold}

The K3\(\times\)E target is a complex threefold:
\(\dim_\C K3\times E=3\).  The integer \(n\) in the
``level-\(n\) shadow'' is
\[
n=\chi(\mathcal O_Y)
\]
for the one-dimensional subscheme \(Y\subset K3\times E\), whereas
\(\chi(\mathcal O_{K3\times E})=2\cdot 0=0\) is the target
holomorphic Euler characteristic.  The perfect
obstruction theories in Chapter~5 are threefold theories; their
virtual classes are reduced virtual classes on the moduli of objects
on the threefold; the Pfaffian \(\mathrm{Pf}_{\mathrm{prot}}\) of
Chapter~5 is the Pfaffian of the threefold's protected operator data.
The Borcherds product formula and the Gritsenko--Nikulin denominator
identity give the Mukai-vector dictionary on the threefold by
projection along the elliptic factor; the OP/DT/PT comparison gives
the scalar shadow on the threefold; the protected Pfaffian gives the
square root.  No fourfold target intervenes anywhere in the chain.

\subsection{Comparison with the determinant-only construction}

The virtual \(K_0\)-determinant identity
\(\mathcal D_X(Z)=\Delta_5(Z)\) is a Grothendieck-class
identity: it sees only the signed superdimensions of the underlying
virtual super vector spaces.  It does not see brackets, not
parities of representatives, not relation constants, not completed
PBW data, not Hopf-pairing radicals, not orientation gerbes, not the
Pfaffian line \(\mathcal L_{\mathrm{Pf}}\).  Squaring the determinant
gives \(\Delta_5^2\); inverting gives the scalar shadow.  All the
information about the retained operator datum \(\mathfrak D_X^{\mathrm{DI}}\)
that distinguishes the retained-data identity
\(\mathrm{Pf}_{\mathrm{prot}}(\mathfrak
D_X^{\mathrm{DI}})=\Delta_5\) from a non-Pfaffian square-root choice
\(\Delta_5\mapsto -\Delta_5\) lives in the Pfaffian datum and the
orientation character.  The inverse-square scalar shadow
\eqref{eq:zbps-equality}, whose OP representative is
\eqref{eq:zbps-op-dt}, is the maximal information visible after the
trace; the Pfaffian and the orientation are the stronger pre-trace data.

\section{Open structures}
\label{sec:zbps-open-problems}

\subsection{The Hall--Borcherds residual revisited}

Conjecture~\ref{conj:zbps-hall-borcherds-residual} asks for the
filtered \(\mathsf{SC}^{\mathrm{ch,top}}\)
morphism from the \(K3\times E\) Hall boundary object to the
gravitational boundary line algebra of Vol~II, whose derived-centre
trace has scalar character \(\Delta_5^{-2}\).  Its trace part is the
bridge from the level-\(\mathsf{S}\) scalar identity
\(Z_{\mathrm{BPS}}^{K3\times E}=\Delta_5^{-2}\) to a
level-\(\mathsf{Z}\) chiral-Hochschild trace datum.  Its acting part is
the further level-\(\mathsf A\) gravity-line operator algebra.

The residual consists of three completions.  The Drinfeld doubling
\(\widehat{\mathrm{CoHA}}^{\mathrm{red}}\to D(\widehat{\mathrm{CoHA}}^{\mathrm{red}})\)
of the Vol~III positive-half pairing: the pairing, completion,
integral form, and stable-envelope transport that convert the
Hall positive half \(\mathsf{Y}^+(K3\times E)\) into the full quantum
group \(G(K3\times E)\).  The current enveloping along \(E\)
and weight-completion: the chiral specialization
\(\mathrm{Sp}^{\mathrm{ch}}_{\Sigma_{d-1},C}\) along the reference
curve \(C=E\), together with the weight-completed ambient required to
host the formal generators of \(\mathfrak g_{\Delta_5}\).  The
filtered \(\mathsf{SC}^{\mathrm{ch,top}}\) morphism property: the
resulting boundary object must be a filtered morphism in the
two-coloured Swiss-cheese chiral-topological operad, not merely a
graded-vector-space comparison.  Their existence is the Vol~II
Hall--Borcherds residual conjecture.

\subsection{The gravity-line operator algebra and the Pentagon-face
determinant}

The Vol~II gravity-line problem asks for a chain-level operator algebra
\(\mathfrak{Op}_{\mathrm{grav}}^{K3\times E}\) whose Pentagon-face
determinant section is \(\Phi_{10}^{\mathrm{un}}=\Delta_5^2\).  If such
an algebra is constructed, its partition character is the reciprocal
section \((\Phi_{10}^{\mathrm{un}})^{-1}=\Delta_5^{-2}\).  The
Brown--Henneaux dictionary \(c=3\ell/(2G_N)\) fixes the Virasoro
boundary datum at \(\mathrm{Vir}_c\); it does not construct the
\(K3\times E\) gravity-line algebra and does not evaluate the
Pentagon-face determinant.  A K3-Heisenberg witness would have to
specialize the Virasoro \(\lambda\)-bracket
\(\{T_\lambda T\}=\partial T+2T\lambda+(c/12)\lambda^3\) to the closed
K3 boundary at \(\kappa_{\mathrm{BKM}}(\Delta_5)=5\) and identify the
resulting determinant with \(\Phi_{10}^{\mathrm{un}}=\Delta_5^2\).
Existence of the gravity-line operator algebra, agreement of its
reciprocal partition character with the unscaled trace
\eqref{eq:zbps-equality}, equivalently with the OP representative
\eqref{eq:zbps-op-dt} after OP normalization on each primitive closed
K3-class fibre, and factorisation through the protected Pfaffian
\(\mathrm{Pf}_{\mathrm{prot}}(\mathfrak D_X^{\mathrm{DI}})=\mathcal D_X=\Delta_5\)
of Chapter~5 are open level-\(\mathsf A\) conditions.

\subsection{Mathieu and umbral moonshine adjacency}

The K3 elliptic genus
\(Z_{\mathrm{K3}}=2\phi_{0,1}\) admits the
\(N=4\) superconformal-character decomposition of
Eguchi--Ooguri--Tachikawa \cite{EOT}: after the massless characters are
removed, the holomorphic mock-modular series
\[
\Sigma(\tau)=-2q^{-1/8}+\sum_{n\ge1}A_nq^{n-1/8}
\]
has coefficients \(A_n\) which are dimensions of, and after twining
characters of, \(M_{24}\)-modules \cite{GHV,GannonUmbral}.  The
decomposition is the infinite massive-sector decomposition of the
\(N=4\) algebra.  The coefficients \(f(n,l)\) of
\(\phi_{0,1}\) which drive the Borcherds product are the untwined
Jacobi coefficients; an \(M_{24}\)-twining would replace
\(\phi_{0,1}\) by the corresponding twined K3 elliptic genus before
the singular theta lift.  The umbral moonshine extension of
Cheng--Duncan--Harvey \cite{CDH} is indexed by the \(23\) Niemeier root
systems with non-empty root system, with \(M_{24}\) the
\(A_1^{24}\) case.  The interaction between this sporadic-group
structure on the K3 elliptic genus and the automorphic / BKM
denominator structure on the
\(K3\times E\) Igusa cusp form is the Mathieu / umbral moonshine
adjacency.

\subsection{The closed-cobordism partition reads as one}

The scalar shadow is the trace, the Pfaffian is the pre-trace
operator-level section, and the primitive input is the retained
\(E_3\)-prefactorization input whose Pfaffian line is presented by
\(\mathfrak D_X^{\mathrm{DI}}\).  The closed
\(K3\times E\to\mathrm{pt}\) cobordism sees only the orientation-forgetful
scalar trace.  At that level the inverse of the squared Pfaffian section is the Laurent expansion
\(\Delta_5^{-2}=(\Phi_{10}^{\mathrm{un}})^{-1}\) of the inverse Igusa
cusp form, and the orientation information that distinguishes
\(\Delta_5\) from \(-\Delta_5\) lives in the Pfaffian datum at level
\(\mathsf{Z}\).  The closed-cobordism
partition function reads as one because \(\Delta_5\) is the same
section throughout: the Borcherds--Gritsenko--Nikulin section of the
weight-\(5\) automorphic line, the Borcherds--Kac--Moody denominator
on the root lattice \(\Lambda^{2,1}_{II}\), and the protected Pfaffian
of the Dirac--Igusa pro-object on \(K3\times E\), relative to the
retained orientation and wall-atlas data.
The squared shadow of this same section is the protected partition
function on the closed cobordism \(K3\times E\to\mathrm{pt}\); the
Hall--Borcherds residual of Vol~II, when constructed, will lift it to
the gravity-line operator algebra at level \(\mathsf{A}\); the
threefold realization question of
\S\ref{sec:k3-e-threefold-realization} is the retained-data
construction of the Dirac--Igusa pro-object and the remaining
compact-source recognition problem.

\section*{Bibliographic notes}

The OP/DT scalar identity is Oberdieck--Pixton \cite{OB},
Oberdieck--Pandharipande \cite{OPand}, and Oberdieck
\cite{OberdieckReducedPT}; the variable conventions follow Bryan
\cite{BRYAN}.  The Oberdieck--Pixton/Pandharipande scalar identity
\cite[Theorem~1]{OB} is the consolidated form used here.  The
Andrianov factorization of the spinor \(L\)-function
is \cite{Andr74}; the Saito--Kurokawa correspondence at weight
\(10\) is \cite{ZagierSK}; the Manin--Deligne period theorem is
\cite{Manin72}, \cite{KZ}, \cite{Del79}.  The dyon-counting
interpretation in CHL backgrounds is Dijkgraaf--Verlinde--Verlinde
\cite{DVV} and Borcherds \cite{B}; these inputs are the
boundary-row physical hosts and are kept separate from the
\(K3\times E\) protected scalar shadow throughout.

The Hall--Borcherds residual conjecture in
\S\ref{sec:hall-borcherds-residual} is Vol~II's Construction
Problem~2 \cite{ChiralBarCobarVolII}, the gravity-line Hall--Borcherds
comparison.
The protected D-brane index conjecture
\(\Omega_X(h,d,n)\stackrel{\mathrm{phys}}{=}\mathbf{DT}^X_{n,(\beta_h,d)}\)
is recorded as the protected D-brane-index conjecture of
Chapter~5; it is not used in the OP/DT scalar proof.  The
Mukai dictionary, the K3-class chamber convention, and the
\(\Lambda^{2,1}_{II}\) lattice data are imported from Chapter~2.  The
protected Pfaffian \(\mathrm{Pf}_{\mathrm{prot}}\) and the Dirac--Igusa
pro-object \(\mathfrak D_X^{\mathrm{DI}}\) are specified in
Chapter~5; their Pfaffian and orientation parts are established
relative to the retained data of Appendix~\ref{app:obstruction-discharges};
the Pfaffian--Dirac theorem is Chapter~6; the four
obstruction classes (D0), (O1), (O1\(^+\)), (O2) are defined and
tracked in \S5.8 of Chapter~5.

The \(\kappa\)-tuple \((0,0,3,5,24)\) and its non-additivity at \(N=1\) (the
controlling instance for \(\Delta_5\)) are recorded in Vol~III; the
Hall positive half \(\mathsf{Y}^+(K3\times E)\) and its Drinfeld
doubling datum toward \(G(K3\times E)\), as well as the 6d holomorphic
Chern--Simons quartic obstruction
\(\int_X\mathrm{Tr}_{\mathrm{ad}}(A(F_A)^3)\), are also Vol~III.  The
local model \(\R^2_{\mathrm{top}}\times\C^2_{\mathrm{hol}}\) and the
holomorphic de Rham obstruction class are mixed-HT-strings.  None of
these supply a 3D gravitational path integral.
